\documentclass[12pt]{article}
\usepackage[utf8]{inputenc}
\usepackage[T1]{fontenc}
\usepackage[margin=1in]{geometry}
\usepackage{setspace}
\usepackage{xcolor}
\usepackage[colorlinks=true, urlcolor=blue!60!black, linkcolor=blue!60!black]{hyperref}
\setlength{\parindent}{0em}
\setlength{\parskip}{0.6em}
\renewcommand{\baselinestretch}{1.0}
\begin{document}
\begin{center}
{\LARGE \textit{Reynolds v. United States}, 98 U.S. 145 (1879)\par}
\vspace{0.5em}
{\large Case Brief --- Two-Page Edition\par}
\end{center}

\section{Memory Jogger}
The Free Exercise Clause protects religious belief, not conduct; a sincere religious duty to practice polygamy is no defense to a federal bigamy law (\href{https://supreme.justia.com/cases/federal/us/98/145/}{Justia}).

\section{Detailed Case Facts}
George Reynolds, a Latter-day Saint in Utah Territory, took a second wife while his first was living, believing polygamy a divinely ordained duty. The Morrill Anti-Bigamy Act of 1862 made bigamy in federal territories a crime. The judge refused to instruct that religious duty required acquittal.

\section{Procedural History}
Convicted in the territorial district court and sentenced to two years at hard labor and a \$500 fine, Reynolds lost in the Supreme Court of Utah Territory and brought a writ of error.

\section{Judicial Votes}
Affirmed, 9--0 on free exercise and 8--1 on confrontation (Waite for the Court; Field dissenting in part). On rehearing, ``at hard labor'' was struck from the sentence.

\section{Holding}
Congress may not regulate religious belief but may punish actions that violate social duties or subvert good order. Religious duty neither exempts a defendant from a general criminal law nor negates intent. A defendant who wrongfully procures a witness's absence forfeits his confrontation objection.

\section{Analysis of Opinions}
\textbf{Waite:} Relying on Jefferson's Danbury letter that government powers ``reach actions only, and not opinions,'' he announced the belief-action line. Polygamy had long been criminal, and marriage underlies social order. Excusing religious human sacrifice would be absurd; so is a religious exemption here. \textbf{Field:} joined on free exercise but found no foundation for admitting the absent witness's former testimony.

\section{Examples: Future Applications}
\textbf{Same-Side:} Sacramental use of an illegal drug may be punished (\textit{Smith}). \textbf{Opposite-Side:} A law punishing the profession of belief in plural marriage is void. \textbf{Fence-Sitter:} Amish parents ending schooling after eighth grade won an exemption in \textit{Yoder}.

\section{Critique}
McConnell faults Waite's reliance on Jefferson over Madison, and critics see anti-Mormon bias in the opinion. Defenders, echoing \textit{Smith}, warn that exemptions threaten the equal rule of law.

\section{Key Quotations}
\begin{quote}
``Can a man excuse his practices to the contrary because of his religious belief? To permit this would be to make the professed doctrines of religious belief superior to the law of the land, and, in effect, to permit every citizen to become a law unto himself. Government could exist only in name under such circumstances.''
\end{quote}
\noindent --- Chief Justice Waite, 98 U.S. at 166--167.

\end{document}
